\documentclass[a4paper]{article}     
\usepackage{amsfonts}
\usepackage{amsmath}
\usepackage{amssymb}
\usepackage{graphicx}

\begin{document}

\noindent \large Auteur: Mateusz Ragan \\
\noindent \large Studierichting: Informatica\\
\noindent \large Oefeningenreeks 2E nr. 26.8\\

\medskip

\normalsize

$\log_{3}(10-3^x) = \log_{3}(81) + \log_{9}\left(\dfrac{1}{81}\right)-x$\\

$\Leftrightarrow \log_{3}(10-3^x) = 4 + \log_{9}(1)-\log_{9}(81)-x$\\

$\Leftrightarrow \log_{3}(10-3^x) = 4 + 0 - 2 - x$\\

$\Leftrightarrow \log_{3}(10-3^x) = 2-x$\\

$\Leftrightarrow 3^{2-x} = 10-3^x$\\

$\Leftrightarrow \dfrac{9}{3^x} = 10-3^x$ stel$3^x = t$\\

$\Leftrightarrow \dfrac{t^2-10t+9}{t} = 0$\\

$\Delta = 100-36 = 64$\\

$t_{1,2} = \dfrac{10\pm8}{2}$\\

$t_{1} = 9$ $t_{2} = 1$\\

$3^x = 9 \Leftrightarrow x = 2$\\

$3^x = 1 \Leftrightarrow x = 0$\\

dus $x = 2 \vee x = 0$




\begin{thebibliography}{999}
%\bibitem{a} Referentie
\end{thebibliography}
\end{document} 
